%%%%%%%%%%%%%%%%%%%%%%%%%%%%%%%%%%%%%%%%%%%%%%%%%%%%%%%%%%%%%%%%%%%%%%%%%%%
% Chapter: Noise and Interference
% Falstad examples: examples/noise (generated by falstad/circuits/noise.py)
% Paper numbers: k T with T = 300 K (4kT = 1.657e-20 J); capstone values as
% in Capstone.tex (bridge 0..37.18 mV, G = 107.6, fc = 2.34 Hz, 3 bit, 4 V).
%%%%%%%%%%%%%%%%%%%%%%%%%%%%%%%%%%%%%%%%%%%%%%%%%%%%%%%%%%%%%%%%%%%%%%%%%%%

% noifig=<s>: scales coordinates and component size by s, but not the text
\tikzset{noifig/.style={scale=#1, /tikz/circuitikz/bipoles/length=#1cm}}
% unit of a noise density, e.g. \SI{4.07}{\nano\volt}\rtHz
\providecommand{\rtHz}{\ensuremath{/\!\sqrt{\mathrm{Hz}}}}

%%%%%%%%%%%%%%%%%%%%%%%%%%%%%%%%%%%%%%%%%%%%%%%%%%%%%%%%%%%%%%%%%%%%%%%%%%%
\section{Signal, Noise and Interference}
%%%%%%%%%%%%%%%%%%%%%%%%%%%%%%%%%%%%%%%%%%%%%%%%%%%%%%%%%%%%%%%%%%%%%%%%%%%
\nsl{Every measured voltage is the wanted signal plus something we do not
want. This chapter collects the unwanted parts that matter for a
measurement chain like the capstone: a sensor signal of a few millivolts,
amplified about a hundred times, filtered and sampled. We distinguish three
groups. \emph{Noise} in the narrow sense is random and is produced inside
the components themselves: thermal noise of every resistor, shot noise of
every current that crosses a junction, and the $1/f$ noise of
semiconductors. It can be reduced but never removed. \emph{Interference}
is a deterministic signal from somewhere else -- the mains, a switching
regulator, the LEDs on the same board -- that reaches our circuit through
a coupling path: an electric field, a magnetic field, or a shared
conductor. It can in principle be removed completely by breaking the path.
\emph{Drift and offset} are slow errors such as thermoelectric voltages,
offset voltages and self-heating; they look like a wrong DC value.
\newline
For each effect we ask three questions: where does it come from, how large
is it in our use case, and what can we do about it. The yardstick is the
capstone: one LSB of the 3-bit converter is \SI{0.5}{\volt} at the ADC,
\SI{4.65}{\milli\volt} at the bridge, or \SI{12.5}{\kelvin}. Because the
same chain with a 12-bit converter would be the obvious next step, we also
compare with its LSB of \SI{0.98}{\milli\volt} at the ADC (\SI{9.1}{\micro\volt}
at the bridge). The previous chapters already treated some disturbances:
the \SI{50}{\hertz} interference in series with the sensor and the lead
resistance (chapter~\ref{ch:Bridges}), offset and bias currents
(chapter~\ref{ch:OpAmps}), the anti-aliasing filter (chapter~\ref{ch:Filters})
and quantization noise (chapter~\ref{ch:ADC}). Here they get their common
framework. Further reading: \cite{horowitz2015} (chapter~8, low-noise
techniques), \cite{ott2009} (grounding, shielding, coupling),
\cite{motchenbacher1993} (noise analysis), \cite{kester2005} and
\cite{tietze2008}.\newpage}

\begin{description}
\item[{\rot\bf Three kinds of unwanted signals}]~\\[-\lblskip]
\begin{center}
\small
\begin{tabular}{>{\raggedright}p{2.2cm}>{\raggedright}p{4.3cm}>{\raggedright}p{4.0cm}>{\raggedright\arraybackslash}p{4.4cm}}
\toprule
 & \textbf{origin} & \textbf{examples} & \textbf{remedy} \\
\midrule
\textbf{noise} & random motion of charges in the components & thermal, shot, $1/f$ noise & less bandwidth, low-noise parts \\
\textbf{interference} & a foreign signal, coupled in through a path & mains hum, switching ripple, LED currents & break the path: layout, shield, differential input, filter \\
\textbf{drift, offset} & slow DC errors & thermal EMF, offset, self-heating & symmetry, calibration, chopping \\
\bottomrule
\end{tabular}
\end{center}
\bi
	\ite noise is random: RMS value and spectrum; interference has a waveform (\SI{50}{\hertz}, a clock)
\ei
\os{\newpage}

\item[{\rot\bf Where disturbances enter the capstone chain}]~\\[-\lblskip]
\begin{center}
\resizebox{\linewidth}{!}{%
\begin{tikzpicture}[font=\sffamily\small, >=Stealth,
  blk/.style={draw, thick, rounded corners=1pt, minimum height=11mm, minimum width=19mm, align=center},
  dis/.style={draw=rwuviolet, fill=rwuvioletlight!15, rounded corners=1pt, font=\sffamily\scriptsize, align=center, inner sep=2pt}]
  \node[blk, fill=rwucyan!10] (b) {Pt100\\bridge};
  \node[blk, fill=rwuviolet!15, right=9mm of b] (i) {INA\\$G=107.6$};
  \node[blk, fill=rwucyan!10, right=9mm of i] (f) {low-pass\\\SI{2.34}{\hertz}};
  \node[blk, fill=rwucyan!10, right=9mm of f] (s) {S\&H\\$f_s=\SI{5}{\hertz}$};
  \node[blk, fill=rwucyan!10, right=9mm of s] (a) {3-bit\\flash ADC};
  \draw[->, thick] (b) -- node[above, font=\scriptsize]{0..37\,mV} (i);
  \draw[->, thick] (i) -- node[above, font=\scriptsize]{0..4\,V} (f);
  \draw[->, thick] (f) -- (s);
  \draw[->, thick] (s) -- (a);
  \node[dis, above=7mm of b] (d1) {50\,Hz fields\\thermal EMF\\lead resistance};
  \node[dis, above=7mm of i] (d2) {op-amp noise\\offset, $1/f$\\RF rectification};
  \node[dis, above=7mm of f] (d3) {resistor noise};
  \node[dis, above=7mm of s] (d4) {aliasing\\$kT/C$};
  \node[dis, above=7mm of a] (d5) {quantization\\reference noise};
  \foreach \n/\m in {d1/b, d2/i, d3/f, d4/s, d5/a} \draw[->, rwuviolet] (\n) -- (\m);
  \node[dis, below=7mm of i, minimum width=7cm] (g) {ground: shared returns (LEDs), ground loops, supply ripple};
  \draw[->, rwuviolet] (g) -| (b); \draw[->, rwuviolet] (g) -| (a);
  \draw[->, rwuviolet] (g) -- (i);
\end{tikzpicture}}
\end{center}
\bi
	\ite the bridge and the amplifier input are the sensitive points: everything there is multiplied by 107.6
	\ite after the low-pass only slow disturbances (below \SI{2.34}{\hertz}) and sampling effects remain
\ei
\os{\newpage}

\item[{\rot\bf The yardstick: one LSB}]~\\[-\lblskip]
\begin{center}
\begin{tabular}{lccc}
\toprule
 & at the ADC & at the bridge & temperature \\
\midrule
full scale & \SI{4.00}{\volt} & \SI{37.18}{\milli\volt} & \SI{100}{\kelvin} \\
1 LSB, 3 bit (capstone) & \SI{0.5}{\volt} & \SI{4.65}{\milli\volt} & \SI{12.5}{\kelvin} \\
1 LSB, 12 bit & \SI{0.98}{\milli\volt} & \SI{9.1}{\micro\volt} & \SI{0.024}{\kelvin} \\
\bottomrule
\end{tabular}
\end{center}
\bi
	\ite bridge sensitivity: \SI{37.18}{\milli\volt}/\SI{100}{\kelvin} $=$ \SI{0.37}{\milli\volt\per\kelvin}
	\ite a disturbance is harmless if it stays well below $\frac12$ LSB (better $\frac14$ LSB)
	\ite the 3-bit capstone tolerates a lot -- with 12 bit almost every effect of this chapter matters
\ei
\nsl{Comparing every disturbance with the LSB at the point where it enters
the chain tells us immediately whether we have to act. A disturbance at
the bridge is compared with \SI{4.65}{\milli\volt} (3 bit) or
\SI{9.1}{\micro\volt} (12 bit); a disturbance at the ADC input with
\SI{0.5}{\volt} or \SI{0.98}{\milli\volt}.}
\os{\newpage}
\end{description}

%%%%%%%%%%%%%%%%%%%%%%%%%%%%%%%%%%%%%%%%%%%%%%%%%%%%%%%%%%%%%%%%%%%%%%%%%%%
\section{Describing Noise}
%%%%%%%%%%%%%%%%%%%%%%%%%%%%%%%%%%%%%%%%%%%%%%%%%%%%%%%%%%%%%%%%%%%%%%%%%%%
\nsl{A noise voltage has no amplitude or phase that we could predict. We
describe it statistically: by its RMS value (the standard deviation of the
voltage), by its amplitude distribution, and by how its power is spread over
frequency, the spectral density. The spectral density is the most useful
description, because it tells us how much noise passes a filter of a given
bandwidth.\newpage}

\begin{description}
\item[{\rot\bf RMS value and peak-to-peak value}]~\\[-\lblskip]
\begin{center}
\begin{tikzpicture}
\begin{axis}[width=12cm, height=3.8cm, xmin=0, xmax=200, ymin=-4, ymax=4,
  xlabel={$t$}, ylabel={$v_n/V_\mathrm{rms}$}, grid=major, xticklabels={},
  ytick={-3.3,0,3.3}, ticklabel style={font=\scriptsize}, label style={font=\small}]
  \pgfmathsetseed{7}
  \addplot[rwuviolet, thin, samples=200, domain=0:200] {1.0*(rand+rand+rand)};
  \addplot[rwucyandark, dashed, thick] coordinates {(0,3.3) (200,3.3)};
  \addplot[rwucyandark, dashed, thick] coordinates {(0,-3.3) (200,-3.3)};
\end{axis}
\end{tikzpicture}
\end{center}
\bi
	\ite most noise is Gaussian: the RMS value $V_\mathrm{rms}$ is the standard deviation $\sigma$
	\ite there is no real peak value; $\pm 3.3\sigma$ contains the signal \SI{99.9}{\percent} of the time
\ei
\keyeq{Peak-to-peak noise (rule of thumb)}{$V_\mathrm{pp} \approx 6.6\cdot V_\mathrm{rms}$}
\nsl{Data sheets often give the noise from \SIrange{0.1}{10}{\hertz} as a
peak-to-peak value, because that is what one sees on an oscilloscope. The
factor 6.6 converts it into an RMS value and back; some manufacturers use
6 instead. For the comparator of a flash converter the peak-to-peak value
is what counts: its hysteresis must be larger than the peak-to-peak noise
(chapter~\ref{ch:ADC}).}
\os{\newpage}

\item[{\rot\bf Spectral density}]~\\[-\lblskip]
\bi
	\ite noise power spreads over frequency; the \textbf{noise density} $e_n$ is the RMS voltage in a band of \SI{1}{\hertz}
	\ite unit: $\mathrm{V}/\sqrt{\mathrm{Hz}}$, e.g.\ \SI{10}{\nano\volt}\rtHz
	\ite \textbf{white} noise: $e_n$ is the same at all frequencies
\ei
\keyeq{White noise in a bandwidth $B$}{$\displaystyle V_\mathrm{rms} = e_n\cdot\sqrt{B} \qquad\qquad \text{in general:}\quad V_\mathrm{rms}^2 = \int e_n^2(f)\,|H(f)|^2\,df$}
\bi
	\ite ten times the bandwidth gives $\sqrt{10} = 3.16$ times the noise
	\ite the narrower the filter, the less noise: \textbf{use only the bandwidth you need}
\ei
\nsl{The power of independent frequency bands adds, so the square of the
density, $e_n^2$ in $\mathrm{V^2/Hz}$, is the quantity that is integrated over
frequency. The square root of the integral is the RMS value. A filter with
transfer function $H(f)$ weights each band with $|H(f)|^2$.}
\os{\newpage}

\item[{\rot\bf Noise bandwidth of a filter}]~\\[-\lblskip]
\bi
	\ite a real filter does not cut off sharply; it passes some noise above $f_c$
	\ite the \textbf{equivalent noise bandwidth} $B_n$ is the width of an ideal (brick-wall) filter that passes the same noise power
\ei
\keyeq{Equivalent noise bandwidth}{$\displaystyle B_n = \int_0^\infty |H(f)|^2\,df$ \qquad
first order: $B_n = \dfrac{\pi}{2}f_c = 1.57\,f_c$ \qquad second-order Butterworth: $B_n = 1.11\,f_0$}
\bi
	\ite capstone low-pass ($f_c = \SI{2.34}{\hertz}$): $B_n = \SI{3.68}{\hertz}$
	\ite white noise behind it: $V_\mathrm{rms} = e_n\sqrt{\SI{3.68}{\hertz}} = 1.92\,e_n\cdot\sqrt{\mathrm{Hz}}$
\ei
\os{\newpage}

\item[{\rot\bf Bandwidth in Falstad}]~\\[-\lblskip]
\bi
	\ite Falstad's noise generator: a new random value at every time step, uniform in $\pm\SI{1}{\volt}$
	\itee RMS $= \SI{1}{\volt}/\sqrt3 = \SI{0.577}{\volt}$, white up to $1/(2\Delta t) = \SI{50}{\kilo\hertz}$ ($\Delta t = \SI{10}{\micro\second}$)
	\itee density $e_n = \SI{0.577}{\volt}\cdot\sqrt{2\Delta t} = \SI{2.58}{\milli\volt}\rtHz$
	\ite behind RC low-passes (scopes show the RMS value):
\ei
\begin{center}
\begin{tabular}{lccc}
\toprule
 & $f_c$ & $B_n$ & RMS \\
\midrule
no filter & -- & \SI{50}{\kilo\hertz} & \SI{0.577}{\volt} \\
\SI{1.6}{\kilo\ohm}, \SI{100}{\nano\farad} & \SI{995}{\hertz} & \SI{1563}{\hertz} & \SI{0.10}{\volt} \\
\SI{16}{\kilo\ohm}, \SI{100}{\nano\farad} & \SI{99.5}{\hertz} & \SI{156}{\hertz} & \SI{0.032}{\volt} \\
\bottomrule
\end{tabular}
\end{center}
\falstad{noise_bandwidth}\par
\nsl{Falstad cannot simulate the thermal noise of its resistors, but its
noise generator produces white noise of a known density, scaled up so
that it is visible. The values in Falstad change a little from run to run
(a few percent), as a measured noise voltage does. One tenth of the
bandwidth gives $1/\sqrt{10}$ of the noise: $0.10/0.032 = 3.2$.}
\os{\newpage}

\item[{\rot\bf Adding noise sources}]~\\[-\lblskip]
\bi
	\ite \textbf{independent} (uncorrelated) sources: the powers add -- root sum of squares (RSS)
\ei
\keyeq{Uncorrelated noise sources}{$\displaystyle V_\mathrm{rms,tot} = \sqrt{V_1^2 + V_2^2 + \ldots}$ \qquad (correlated: $V_1 + V_2$)}
\bi
	\ite two equal sources: $\sqrt2$ times one source, not twice
	\ite one source three times smaller than the largest adds only \SI{5}{\percent}: $\sqrt{1 + (1/3)^2} = 1.054$
	\itee \textbf{only the largest contributor matters} -- reduce that one first
	\ite Falstad: one generator \SI{0.577}{\volt}; two independent ones in series \SI{0.82}{\volt}; one generator amplified by 2 (fully correlated) \SI{1.15}{\volt}
\ei
\falstad{noise_rss}\par
\nsl{The rule holds because the cross term $2\,\overline{v_1 v_2}$ of the
mean square $\overline{(v_1+v_2)^2}$ is zero for independent random
signals. The thermal noise of two resistors, the voltage and the current
noise of an op-amp, or the noise of two different stages are independent.
Interference is not noise: two \SI{50}{\hertz} contributions from the same
mains add with their phases.}
\os{\newpage}

\item[{\rot\bf Signal-to-noise ratio}]~\\[-\lblskip]
\keyeq{Signal-to-noise ratio}{$\displaystyle \mathrm{SNR} = 20\log_{10}\frac{V_\mathrm{signal,rms}}{V_\mathrm{noise,rms}}\,\mathrm{dB}$}
\bi
	\ite an ideal $N$-bit ADC with a full-scale sine: $\mathrm{SNR} = 6.02\,N + 1.76\dB$ (chapter~\ref{ch:ADC})
	\itee its quantization noise is $\LSB/\sqrt{12}$: capstone \SI{144}{\milli\volt}, 12 bit \SI{0.28}{\milli\volt}
	\ite the analog noise should be smaller than the quantization noise -- otherwise the extra bits are wasted
	\ite for a slow DC measurement (temperature) we compare the peak-to-peak noise with the LSB
\ei
\os{\newpage}
\end{description}

%%%%%%%%%%%%%%%%%%%%%%%%%%%%%%%%%%%%%%%%%%%%%%%%%%%%%%%%%%%%%%%%%%%%%%%%%%%
\section{Intrinsic Noise Sources}
%%%%%%%%%%%%%%%%%%%%%%%%%%%%%%%%%%%%%%%%%%%%%%%%%%%%%%%%%%%%%%%%%%%%%%%%%%%
\nsl{Three physical mechanisms produce the noise of electronic components.
Thermal noise comes from the random thermal motion of the charge carriers
in every resistance. Shot noise comes from the discreteness of the
charge: a current across a barrier is a stream of independent electrons.
$1/f$ noise comes from slow fluctuations of the number or the mobility of
carriers, for example by traps in the crystal or at surfaces. The first two
are fundamental and can be calculated exactly; the third is a property of
the technology and must be taken from the data sheet.\newpage}

\begin{description}
\item[{\rot\bf Thermal noise (Johnson-Nyquist noise)}]~\\[-\lblskip]
\bi
	\ite origin: the electrons in a conductor move randomly with the thermal energy $kT$ (Brownian motion)
	\itee at any instant slightly more electrons are at one end: a fluctuating open-circuit voltage
	\ite white up to far above \SI{1}{\tera\hertz}; Gaussian; independent of the material and of a DC current
	\ite measured by Johnson, explained by Nyquist in 1928 \cite{johnson1928,nyquist1928}
\ei
\keyeq{Thermal noise of a resistance $R$}{$\displaystyle e_n = \sqrt{4kTR} \qquad V_\mathrm{rms} = \sqrt{4kTRB}
\qquad 4kT = \SI{1.66e-20}{\joule} \text{ at } \SI{300}{\kelvin}$}
\bi
	\ite model: a noise voltage source in series with an ideal $R$ (or a current source $\sqrt{4kT/R}$ in parallel)
	\ite only the real part of an impedance makes noise: capacitors and inductors are noiseless
\ei
\os{\newpage}

\item[{\rot\bf How large is thermal noise?}]~\\[-\lblskip]
\begin{center}
\begin{tabular}{rccc}
\toprule
$R$ & $e_n$ & in \SI{3.68}{\hertz} (capstone) & in \SI{20}{\kilo\hertz} (audio) \\
\midrule
\SI{100}{\ohm} & \SI{1.29}{\nano\volt}\rtHz & \SI{2.5}{\nano\volt} & \SI{0.18}{\micro\volt} \\
\SI{1}{\kilo\ohm} & \SI{4.07}{\nano\volt}\rtHz & \SI{7.8}{\nano\volt} & \SI{0.58}{\micro\volt} \\
\SI{10}{\kilo\ohm} & \SI{12.9}{\nano\volt}\rtHz & \SI{25}{\nano\volt} & \SI{1.8}{\micro\volt} \\
\SI{1}{\mega\ohm} & \SI{129}{\nano\volt}\rtHz & \SI{0.25}{\micro\volt} & \SI{18}{\micro\volt} \\
\bottomrule
\end{tabular}
\end{center}
\bi
	\ite memorize: \textbf{\SI{1}{\kilo\ohm} $\approx$ \SI{4}{\nano\volt}\rtHz}, scale with $\sqrt{R}$
	\ite capstone bridge: each tap sees $\SI{3.9}{\kilo\ohm}\parallelto\SI{100}{\ohm} = \SI{97.5}{\ohm}$; both taps \SI{195}{\ohm}: \SI{1.8}{\nano\volt}\rtHz
	\itee in $B_n = \SI{3.68}{\hertz}$: \SI{3.4}{\nano\volt} RMS -- 1/2700 of a 12-bit LSB (\SI{9.1}{\micro\volt})
\ei
\nsl{Thermal noise sets the lower limit for any measurement with a
resistive source. For the Pt100 bridge this limit is far away: even a
20-bit converter would have an LSB of \SI{35}{\nano\volt} at the bridge,
ten times the bridge noise. Thermal noise becomes the limit for
high-ohmic sources (pH electrodes, piezo sensors, photodiode
transimpedance amplifiers) and for wide bandwidths (audio, radio).
Cooling helps only with the square root of the absolute temperature;
reducing $R$ and $B$ helps much more.}
\os{\newpage}

\item[{\rot\bf $kT/C$ noise: sampling freezes thermal noise}]~\\[-\lblskip]
\begin{center}
\begin{circuitikz}[noifig=1.1]
  \draw (0,0) to[sV, l=$\sqrt{4kTR}$] (0,2) to[R, l=$R$ ($R_\mathrm{on}$)] (3,2) -- (4,2) to[C, l=$C$] (4,0) -- (0,0);
  \draw (4,2) to[short, -o] (5.2,2) node[right]{$v_C$};
  \draw (4,0) node[ground]{};
\end{circuitikz}
\end{center}
\bi
	\ite the resistor noise is filtered by $C$: $e_n^2 \cdot B_n = 4kTR\cdot\dfrac{\pi}{2}\cdot\dfrac{1}{2\pi RC} = \dfrac{kT}{C}$ -- \textbf{independent of $R$}
	\ite when the switch of a sample and hold opens, this noise voltage is held
\ei
\keyeq{$kT/C$ noise}{$V_\mathrm{rms} = \sqrt{kT/C}$: \quad \SI{1}{\pico\farad}: \SI{64}{\micro\volt}, \quad \SI{10}{\pico\farad}: \SI{20}{\micro\volt}, \quad \SI{1}{\micro\farad}: \SI{64}{\nano\volt}}
\bi
	\ite capstone $C_H = \SI{1}{\micro\farad}$ at the 4-V level: negligible; on-chip S\&H capacitors of a few \si{\pico\farad} limit the resolution of fast ADCs
\ei
\os{\newpage}

\item[{\rot\bf Shot noise}]~\\[-\lblskip]
\bi
	\ite origin: a current across a potential barrier (pn junction, vacuum, tunnel barrier) consists of single electrons that cross \textbf{independently} at random times
	\itee like rain on a roof: the mean flow is constant, the number per time interval fluctuates
	\ite white, Gaussian, proportional to $\sqrt{I}$
\ei
\keyeq{Shot noise of a DC current $I$}{$\displaystyle i_n = \sqrt{2qI} \qquad I_\mathrm{rms} = \sqrt{2qIB} \qquad q = \SI{1.602e-19}{\coulomb}$}
\bi
	\ite \SI{1}{\milli\ampere}: \SI{17.9}{\pico\ampere}\rtHz; \quad \SI{1}{\micro\ampere}: \SI{0.57}{\pico\ampere}\rtHz; \quad \SI{1}{\nano\ampere}: \SI{18}{\femto\ampere}\rtHz
	\ite relative noise $i_n/I = \sqrt{2q/I}$ grows for \textbf{small currents}: photodiodes, op-amp bias currents
	\ite a current through a metal resistor from a voltage source has \textbf{no} shot noise (the electrons are not independent), only thermal noise
\ei
\os{\newpage}

\item[{\rot\bf Shot noise in our use case}]~\\[-\lblskip]
\bi
	\ite sensor current \SI{1}{\milli\ampere}
	\itee from the \SI{3.9}{\kilo\ohm} resistor of the bridge: no shot noise
	\itee from a transistor current source (chapter~\ref{ch:Sources}): \SI{17.9}{\pico\ampere}\rtHz $\times$ \SI{100}{\ohm} = \SI{1.8}{\nano\volt}\rtHz -- as much as the bridge resistors, still negligible
	\ite op-amp input current noise $i_n$ (shot noise of the bias current, e.g.\ $I_B = \SI{20}{\nano\ampere}$: \SI{80}{\femto\ampere}\rtHz)
	\itee flows through the source resistance: $i_n R_S = \SI{80}{\femto\ampere}\rtHz \times \SI{97.5}{\ohm} = \SI{8}{\pico\volt}\rtHz$ -- nothing
	\itee but through \SI{1}{\mega\ohm}: \SI{80}{\nano\volt}\rtHz, more than the op-amp's voltage noise
	\ite Zener diodes above about \SI{6}{\volt} (avalanche): much more noise than shot noise -- reason for bandgap references (chapter~\ref{ch:References})
\ei
\os{\newpage}

\item[{\rot\bf $1/f$ noise (flicker noise)}]~\\[-\lblskip]
\begin{center}
\begin{tikzpicture}
\begin{loglogaxis}[width=12cm, height=5.2cm, xmin=0.01, xmax=1e4, ymin=10, ymax=10000,
  xlabel={$f$ in \si{\hertz}}, ylabel={$e_n$ in \si{\nano\volt}\rtHz}, grid=major,
  ticklabel style={font=\scriptsize}, label style={font=\small},
  legend style={font=\scriptsize, at={(0.98,0.95)}, anchor=north east}]
  \addplot[draw=none, fill=rwucyan!20, forget plot] coordinates {(0.01,10) (0.01,10000) (2.34,10000) (2.34,10)} -- cycle;
  \addplot[rwugray, dashed, forget plot] coordinates {(100,10) (100,10000)};
  \addplot[rwuviolet, very thick, domain=-2:4, samples=100] ({10^x}, {40*sqrt(1+100/10^x)});
  \addplot[rwucyandark, very thick, dashed, domain=-2:4, samples=2] ({10^x}, {55});
  \legend{bipolar{,} $e_w = 40$\,nV/$\sqrt{\mathrm{Hz}}${,} $f_{nc} = \SI{100}{\hertz}$, zero-drift (chopper){,} 55\,nV/$\sqrt{\mathrm{Hz}}$}
\end{loglogaxis}
\end{tikzpicture}
\end{center}
\bi
	\ite below the \textbf{corner frequency} $f_{nc}$ the density rises as $1/\sqrt{f}$ (power $\propto 1/f$)
	\ite origin: carriers trapped and released by defects (surfaces, oxide, grain boundaries): slow fluctuations of the conductance
	\ite shaded: the band of the capstone (\SI{0.01}{} \ldots \SI{2.34}{\hertz}) lies completely in the $1/f$ region
\ei
\os{\newpage}

\item[{\rot\bf $1/f$ noise: every decade counts the same}]~\\[-\lblskip]
\keyeq{Noise with white floor $e_w$ and corner $f_{nc}$, from $f_1$ to $f_2$}{$\displaystyle e_n^2(f) = e_w^2\left(1 + \frac{f_{nc}}{f}\right) \qquad
V_\mathrm{rms} = e_w\sqrt{f_{nc}\ln\frac{f_2}{f_1} + (f_2 - f_1)}$}
\bi
	\ite each decade of frequency contributes the same $1/f$ noise power: $f_{nc}\ln 10$
	\ite no lower limit at DC: $f_1 \approx 1/T_\mathrm{meas}$ (one observation of \SI{100}{\second}: $f_1 = \SI{0.01}{\hertz}$)
	\ite example: $e_w = \SI{40}{\nano\volt}\rtHz$, $f_{nc} = \SI{100}{\hertz}$, \SIrange{0.01}{3.68}{\hertz}
	\itee $1/f$ part: $100\cdot\ln 368 = 591$, white part: $3.67$ (both in \si{\hertz})
	\itee $V_\mathrm{rms} = \SI{40}{\nano\volt}\rtHz\cdot\sqrt{\SI{595}{\hertz}} = \SI{0.98}{\micro\volt}$ -- 13 times the white-noise estimate
\ei
\nsl{Data sheets therefore specify the low-frequency noise directly as a
peak-to-peak voltage from \SIrange{0.1}{10}{\hertz}, for example
\SI{0.35}{\micro\volt}\,pp for the precision bipolar op-amp OP07 or
\SI{1.1}{\micro\volt}\,pp for the zero-drift op-amp OPA333. The $1/f$
corner of general-purpose CMOS op-amps can lie at \SI{1}{\kilo\hertz} or
higher. Resistors show $1/f$ noise as well, as \emph{excess noise} that
grows with the DC voltage across them: thick-film and carbon resistors
are much worse than metal-film and thin-film resistors. In a DC
measurement path use metal-film resistors. A related effect is
\emph{popcorn} or burst noise: random steps of the offset between two or
more levels, caused by single defects; it occurs in poor processes and is
screened out by good manufacturers.}
\os{\newpage}

\item[{\rot\bf Remedies against $1/f$ noise}]~\\[-\lblskip]
\bi
	\ite \textbf{zero-drift op-amps} (chopper, auto-zero): the input is modulated to a high frequency, amplified there and demodulated
	\itee no $1/f$ noise, offset $<\SI{10}{\micro\volt}$, drift $<\SI{0.05}{\micro\volt\per\kelvin}$
	\itee price: higher white noise (e.g.\ \SI{55}{\nano\volt}\rtHz) and small spikes at the chopper frequency
	\ite \textbf{AC excitation}: excite the bridge with a square wave and demodulate synchronously -- the signal moves above $f_{nc}$ (same idea as the chopper, applied to the whole chain)
	\ite metal-film resistors in the signal path
	\ite calibration does not help: $1/f$ noise is a random drift
\ei
\os{\newpage}

\item[{\rot\bf Noise model of the op-amp}]~\\[-\lblskip]
\begin{center}
\begin{circuitikz}[noifig=1.05]
  \draw (0,0) node[op amp, noinv input up](OA){};
  \coordinate (X) at ($(OA.out)+(0.5,0)$);
  \draw (OA.+) -- ++(-0.4,0) coordinate(p) to[sV, l_=$e_n$] ++(-1.8,0) coordinate(q) to[R, l_=$R_S$] ++(-2.0,0) -- ++(0,-0.7) node[ground]{};
  \draw (q) to[sI, l_=$i_n$] ++(0,-1.6) node[ground]{};
  \draw (q) node[circ]{};
  \draw (OA.-) -- ++(0,-0.8) coordinate(N);
  \draw (N) to[R, l_=$R_2$] (N -| X) -- (X);
  \draw (N) to[R, l_=$R_1$] ++(0,-1.4) node[ground]{};
  \draw (OA.out) -- (X) to[short,-o] ++(0.8,0) node[right]{$\Vout$};
\end{circuitikz}
\end{center}
\bi
	\ite like offset and bias current: a noise voltage $e_n$ and noise currents $i_n$ at the inputs
	\ite input-referred noise of a non-inverting amplifier (all terms independent):
\ei
\keyeq{Input noise, non-inverting amplifier}{$e_\mathrm{in}^2 = e_n^2 + 4kTR_S + (i_nR_S)^2 + 4kT\,(R_1\parallelto R_2) + \left(i_n\,(R_1\parallelto R_2)\right)^2$\par
at the output: $\left(1 + R_2/R_1\right) e_\mathrm{in}$ \quad (noise gain)}
\os{\newpage}

\item[{\rot\bf Choosing the op-amp by its noise}]~\\[-\lblskip]
\begin{center}
\begin{tabular}{lccc}
\toprule
type (example) & $e_n$ at \SI{1}{\kilo\hertz} & $i_n$ & $1/f$ \\
\midrule
general purpose bipolar (LM358) & \SI{40}{\nano\volt}\rtHz & medium & yes \\
precision bipolar (OP07) & \SI{10}{\nano\volt}\rtHz & \SI{0.1}{\pico\ampere}\rtHz & yes, \SI{0.35}{\micro\volt}\,pp \\
CMOS (MCP6002) & \SI{30}{}\ldots\SI{60}{\nano\volt}\rtHz & \si{\femto\ampere}\rtHz & yes, high corner \\
zero-drift (OPA333) & \SI{55}{\nano\volt}\rtHz & \SI{0.1}{\pico\ampere}\rtHz & none, \SI{1.1}{\micro\volt}\,pp \\
\bottomrule
\end{tabular}
\end{center}
\bi
	\ite the source resistance decides: noise-optimal at $R_S \approx e_n/i_n$
	\itee low $R_S$ (bridge, \SI{100}{\ohm}): low $e_n$ matters (bipolar); high $R_S$ (\si{\mega\ohm}): low $i_n$ matters (CMOS, JFET)
	\ite DC measurement: the $1/f$ noise and the offset drift matter most -- zero-drift
\ei
\nsl{The values are typical and rounded; always take the data sheet of the
part that is actually used. For a source resistance $R_S$ the three input
terms are $e_n$, $\sqrt{4kTR_S}$ and $i_nR_S$. With small $R_S$ the voltage
noise dominates, with large $R_S$ the current noise. The resistor noise
$\sqrt{4kTR_S}$ grows only with $\sqrt{R_S}$; an amplifier is ``good
enough'' when its own noise is below the thermal noise of the source.}
\os{\newpage}

\item[{\rot\bf Noise budget of the capstone front end}]~\\[-\lblskip]
\bi
	\ite INA with zero-drift op-amps (\SI{55}{\nano\volt}\rtHz each), \SI{10}{\kilo\ohm}/\SI{100}{\kilo\ohm} difference amplifier
\ei
\begin{center}
\begin{tabular}{lcc}
\toprule
contribution (input-referred) & density & RMS in \SI{3.68}{\hertz} \\
\midrule
bridge resistors (\SI{195}{\ohm}) & \SI{1.8}{\nano\volt}\rtHz & \SI{3.4}{\nano\volt} \\
two input buffers, $\sqrt{2}\cdot\SI{55}{}$ & \SI{78}{\nano\volt}\rtHz & \SI{0.15}{\micro\volt} \\
difference amplifier op-amp, $1.1\cdot\SI{55}{}$ & \SI{61}{\nano\volt}\rtHz & \SI{0.12}{\micro\volt} \\
difference amplifier resistors & \SI{19}{\nano\volt}\rtHz & \SI{0.04}{\micro\volt} \\
\midrule
total (RSS) & \SI{100}{\nano\volt}\rtHz & \SI{0.19}{\micro\volt} \\
\bottomrule
\end{tabular}
\end{center}
\bi
	\ite at the output: $\times 107.6$ = \SI{21}{\micro\volt} RMS, about \SI{0.14}{\milli\volt}\,pp
	\ite 3 bit: 1/3600 LSB; 12 bit: 0.14 LSB\,pp -- the noise of the op-amps is \textbf{not} our problem
\ei
\nsl{The op-amps dominate, the bridge resistors contribute almost nothing.
The second stage (gain 10.76) adds its noise divided by the gain of the
first stage, which is why the first stage should have a gain of 10 or
more. The anti-aliasing filter defines the bandwidth. Without it the
noise bandwidth would be set by the op-amps (some \SI{10}{\kilo\hertz}),
and the noise would be 50 times larger -- and aliased into the sampled
data. With a general-purpose bipolar op-amp ($1/f$ corner
\SI{100}{\hertz}) each op-amp contributes about \SI{1}{\micro\volt}, and
the output noise is about \SI{0.19}{\milli\volt} RMS or \SI{1.2}{\milli\volt}
peak-to-peak (exercise~\ref{ex:noise-ina}): about one LSB of a 12-bit
converter. For the 3-bit capstone the
quantization noise (\SI{144}{\milli\volt} RMS) dominates everything by
orders of magnitude.}
\os{\newpage}
\end{description}

%%%%%%%%%%%%%%%%%%%%%%%%%%%%%%%%%%%%%%%%%%%%%%%%%%%%%%%%%%%%%%%%%%%%%%%%%%%
\section{Coupling of Interference}
%%%%%%%%%%%%%%%%%%%%%%%%%%%%%%%%%%%%%%%%%%%%%%%%%%%%%%%%%%%%%%%%%%%%%%%%%%%
\nsl{Interference needs three things: a source (the mains, a motor, a
switching regulator, a clock line), a victim (our high-gain input), and a
coupling path between them. We can rarely switch off the source and we
cannot make the victim insensitive, so the work is on the path. There are
four paths: a common impedance (a conductor shared by both circuits), an
electric field (capacitive coupling), a magnetic field (inductive coupling),
and electromagnetic radiation. In a laboratory the dominant source is the
mains: 230\,V, 50\,Hz, with harmonics, in every wall and every lamp
\cite{ott2009}.\newpage}

\begin{description}
\item[{\rot\bf Source, path, victim}]~\\[-\lblskip]
\begin{center}
\begin{tikzpicture}[font=\sffamily\small, >=Stealth,
  blk/.style={draw, thick, rounded corners=1pt, minimum height=12mm, minimum width=27mm, align=center}]
  \node[blk, fill=rwuvioletlight!20, minimum height=30mm] (s) {source\\mains, motor,\\regulator, LEDs};
  \node[blk, fill=rwucyan!10, right=75mm of s, minimum height=30mm] (v) {victim\\sensor, INA input,\\reference};
  \foreach \y/\t/\st in {1.05/common impedance (shared conductor)/solid, 0.35/electric field (capacitive)/solid,
      -0.35/magnetic field (inductive)/solid, -1.05/radiation (radio frequencies)/dashed}
    \draw[->, thick, \st] ($(s.east)+(0,\y)$) -- node[above, font=\small]{\t} ($(v.west)+(0,\y)$);
\end{tikzpicture}
\end{center}
\bi
	\ite electric field: from a \textbf{voltage} (mains wiring, switching node): couples into \textbf{high-impedance} nodes
	\ite magnetic field: from a \textbf{current} (mains cable, transformer, motor): couples into \textbf{loops}
	\ite common impedance: a \textbf{current} of one circuit flows through a conductor of the other
\ei
\os{\newpage}

\item[{\rot\bf Capacitive coupling}]~\\[-\lblskip]
\begin{center}
\begin{circuitikz}[noifig=1.1]
  \draw (0,0) to[sV, l=\SI{230}{\volt}] (0,2.2) -- (1,2.2) to[C, l=$C_S$ (stray)] (3,2.2) -- (4,2.2);
  \draw (4,2.2) to[R, l=$Z_V$] (4,0) -- (0,0);
  \draw (4,2.2) to[short, -o] (5.3,2.2) node[right]{$v_V$};
  \draw (2,0) node[ground]{};
\end{circuitikz}
\end{center}
\bi
	\ite the stray capacitance $C_S$ (a few \si{\pico\farad} between cables) injects a current $i = C_S\,dV_S/dt$
	\ite for $Z_V \ll 1/(\omega C_S)$ the victim voltage is this current times the node impedance:
\ei
\keyeq{Capacitive coupling}{$\hat v_V \approx \omega C_S\,\hat V_S\,Z_V$ \qquad mains, \SI{1}{\pico\farad}: $\hat\imath = 2\pi\cdot\SI{50}{\hertz}\cdot\SI{1}{\pico\farad}\cdot\SI{325}{\volt} = \SI{0.10}{\micro\ampere}$}
\os{\newpage}

\item[{\rot\bf Capacitive coupling in Falstad}]~\\[-\lblskip]
\begin{center}
\begin{tabular}{lcc}
\toprule
victim node (\SI{1}{\pico\farad} to the mains) & amplitude & after $G = 107.6$ \\
\midrule
\SI{1}{\mega\ohm} (open input, high-ohmic divider) & \SI{0.10}{\volt} & -- (clipped) \\
\SI{10}{\kilo\ohm} & \SI{1.0}{\milli\volt} & \SI{0.11}{\volt} \\
\SI{100}{\ohm} (bridge tap: \SI{97.5}{\ohm}) & \SI{10}{\micro\volt} & \SI{1.1}{\milli\volt} \\
\SI{1}{\mega\ohm}, shielded (\SI{0.01}{\pico\farad} remain) & \SI{1.0}{\milli\volt} & \SI{0.11}{\volt} \\
\bottomrule
\end{tabular}
\end{center}
\falstad{noise_cap_coupling}\par
\bi
	\ite remedies: \textbf{low node impedance}; distance (smaller $C_S$); a grounded \textbf{shield} takes the field lines
	\ite the low-ohmic bridge is robust; an open or high-ohmic input is not (touch a scope probe!)
	\ite coupled into both bridge wires equally: common mode, rejected by the INA
\ei
\os{\newpage}

\item[{\rot\bf Shielding against electric fields}]~\\[-\lblskip]
\bi
	\ite a conductive shield around the cable, connected to the reference potential, catches the field: the current flows to ground, not into the victim
	\ite connect the shield to the \textbf{signal ground at the amplifier} -- at \textbf{one end} only (low frequencies, see ground loops)
	\ite never let the shield carry signal current (no shield as return conductor)
	\ite cable capacitance of the shield (\SI{100}{\pico\farad\per\metre}) loads high-ohmic sources
	\ite a shield does \textbf{nothing} against \SI{50}{\hertz} magnetic fields (copper is not magnetic)
	\ite guard: a conductor at the same potential as the input (driven by a buffer) around high-impedance nodes -- no leakage current, no field
\ei
\os{\newpage}

\item[{\rot\bf Inductive (magnetic) coupling}]~\\[-\lblskip]
\begin{center}
\begin{tikzpicture}[>=Stealth, font=\small]
  \draw[thick] (0,0) rectangle (5,1.2);
  \node[font=\footnotesize] at (2.5,0.6) {loop area $A$};
  \draw[thick, fill=white] (-0.25,0.35) rectangle (0.25,0.85); \node[left, font=\footnotesize] at (-0.3,0.6) {sensor};
  \draw[thick, fill=rwucyan!20] (4.75,0.35) rectangle (5.25,0.85); \node[right, font=\footnotesize] at (5.3,0.6) {amplifier};
  \foreach \x in {0.6,1.4,...,4.6} \node[rwuviolet] at (\x,1.6) {$\odot$};
  \node[rwuviolet, right, font=\footnotesize] at (5.3,1.6) {$B$ from a mains current};
  \foreach \x in {7.5,8.5,...,12} {\draw[thick] (\x,0.2) .. controls (\x+0.25,0.2) and (\x+0.25,1.0) .. (\x+0.5,1.0) .. controls (\x+0.75,1.0) and (\x+0.75,0.2) .. (\x+1,0.2);
    \draw[thick] (\x,1.0) .. controls (\x+0.25,1.0) and (\x+0.25,0.2) .. (\x+0.5,0.2) .. controls (\x+0.75,0.2) and (\x+0.75,1.0) .. (\x+1,1.0);}
  \node[font=\footnotesize] at (10.25,-0.2) {twisted pair: induced voltages alternate in sign};
\end{tikzpicture}
\end{center}
\keyeq{Voltage induced in a loop}{$\displaystyle \hat v = \omega\,A\,\hat B \qquad \text{field of a long conductor: } B = \frac{\mu_0 I}{2\pi r}$}
\bi
	\ite \SI{10}{\ampere} (RMS) in a mains cable at \SI{10}{\centi\metre}: $\hat B = \SI{28}{\micro\tesla}$
	\ite sensor cable \SI{2}{\metre}, wires \SI{5}{\milli\metre} apart: $A = \SI{100}{\centi\metre\squared}$: $\hat v = \SI{89}{\micro\volt}$
	\ite near a transformer or a lamp ballast ($\hat B \approx \SI{1.6}{\milli\tesla}$): \SI{5}{\milli\volt} -- the capstone value
\ei
\os{\newpage}

\item[{\rot\bf Magnetic coupling: differential mode}]~\\[-\lblskip]
\bi
	\ite the induced voltage is in \textbf{series} with the sensor: a differential signal, amplified like the measurement
	\itee capstone: \SI{5}{\milli\volt} $\rightarrow$ \SI{4.88}{\milli\volt} at the tap $\rightarrow$ \SI{0.52}{\volt} at the INA output (chapter~\ref{ch:Bridges})
	\ite the INA cannot reject it; a shield does not stop it
	\ite remedies
	\itee \textbf{small loop area}: route both wires together, \textbf{twisted pair} -- neighbouring twists cancel (\SIrange{20}{40}{\deci\bel})
	\itee distance from transformers, motors, mains cables; perpendicular crossing
	\itee what remains: \textbf{low-pass filter} (chapter~\ref{ch:Filters}): $-\SI{26.6}{\deci\bel}$ at \SI{50}{\hertz}, \SI{25}{\milli\volt} left
	\ite Falstad: wide loop \SI{5}{\milli\volt}, twisted pair \SI{2.5}{\milli\volt} $-$ \SI{2.25}{\milli\volt} (\SI{10}{\percent} imbalance): \SI{48.8}{\milli\volt} versus \SI{2.4}{\milli\volt} behind the INA ($G = 10$)
\ei
\falstad{noise_mag_coupling}\quad\falstad{bridge_50hz}\par
\os{\newpage}

\item[{\rot\bf Common mode and differential mode}]~\\[-\lblskip]
\bi
	\ite \textbf{common mode} (CM): the same disturbance on both wires against ground (electric field on a symmetric cable, ground potential differences)
	\itee rejected by a differential input: INA with CMRR \SIrange{80}{120}{\deci\bel}
	\ite \textbf{differential mode} (DM): the disturbance between the wires (magnetic field in the loop)
	\itee looks exactly like the signal: avoid it (twisting) or filter it
	\ite \textbf{CM-to-DM conversion}: unequal impedances on the two wires turn common mode into differential mode
	\itee e.g.\ different source resistances with the cable capacitances: keep the two inputs symmetric
\ei
\os{\newpage}

\item[{\rot\bf Radio frequencies and rectification}]~\\[-\lblskip]
\begin{center}
\begin{circuitikz}[noifig=0.9]
  \draw (0,0) node[op amp](OA){};
  \draw (OA.-) -- ++(-0.5,0) |- ++(-0.5,0.9) coordinate(a) -- ++(-1.6,0) coordinate(c1) to[R, l_=\SI{1}{\kilo\ohm}, -o] ++(-2.4,0) node[left]{$V_S$};
  \draw (OA.+) -- ++(-0.5,0) |- ++(-0.5,-0.9) coordinate(b) -- ++(-1.6,0) coordinate(c2) to[R, l=\SI{1}{\kilo\ohm}, -o] ++(-2.4,0) node[left]{$V_R$};
  \draw (a) to[C, l=$C_\mathrm{D}$, *-*] (b);
  \draw (c1) to[C, l_=$C_\mathrm{CM}$, *-] ++(0,1.4) node[ground, rotate=180]{};
  \draw (c2) to[C, l=$C_\mathrm{CM}$, *-] ++(0,-1.4) node[ground]{};
  \draw (OA.out) to[short, -o] ++(0.6,0) node[right]{INA};
\end{circuitikz}
\end{center}
\bi
	\ite mobile phones, WLAN, switching edges: \si{\mega\hertz} to \si{\giga\hertz}, far above the op-amp bandwidth
	\ite the nonlinear input stage \textbf{rectifies} them: a DC shift appears that no later filter removes
	\ite remedy: RC filter \textbf{at the input}, e.g.\ \SI{1}{\kilo\ohm}, $C_\mathrm{D} = \SI{1}{\micro\farad}$, $C_\mathrm{CM} = \SI{10}{\nano\farad}$ (match the CM capacitors, $C_\mathrm{D} \geq 10\,C_\mathrm{CM}$)
\ei
\nsl{The differential capacitor $C_\mathrm{D}$ sets the differential cutoff
frequency, $1/(2\pi\cdot 2R\,(C_\mathrm{D} + C_\mathrm{CM}/2)) \approx
\SI{80}{\hertz}$; the two common-mode capacitors give
$1/(2\pi R C_\mathrm{CM}) = \SI{16}{\kilo\hertz}$ for common-mode signals.
A mismatch of the two common-mode capacitors converts common mode into
differential mode, so $C_\mathrm{D}$ is made at least ten times larger: it
shorts the difference. Falstad cannot show the rectification, because its
op-amp is an ideal linear amplifier \cite{kitchin2006}.}
\os{\newpage}
\end{description}

%%%%%%%%%%%%%%%%%%%%%%%%%%%%%%%%%%%%%%%%%%%%%%%%%%%%%%%%%%%%%%%%%%%%%%%%%%%
\section{Grounding and Ground Loops}
%%%%%%%%%%%%%%%%%%%%%%%%%%%%%%%%%%%%%%%%%%%%%%%%%%%%%%%%%%%%%%%%%%%%%%%%%%%
\nsl{In a schematic all ground symbols are the same node. On a real board
ground is a copper trace or plane with resistance and inductance, and
every current that returns through it causes a voltage drop. Two ground
points of the same circuit are therefore at slightly different potentials,
and the difference depends on what the other circuits are doing at the
moment. The voltages are small -- millivolts -- but at the input of an
amplifier with a gain of 100 they are as large as our signal. This section
treats the two typical cases: a common return impedance on the board, and a
ground loop between two devices.\newpage}

\begin{description}
\item[{\rot\bf Ground is not 0\,V everywhere}]~\\[-\lblskip]
\bi
	\ite copper: $\rho = \SI{0.0172}{\ohm\milli\metre\squared\per\metre}$
	\itee PCB trace \SI{0.3}{\milli\metre} wide, \SI{35}{\micro\metre} thick, \SI{10}{\centi\metre} long: \SI{0.16}{\ohm}
	\itee cable \SI{0.14}{\milli\metre\squared}, \SI{2}{\metre}: \SI{0.25}{\ohm}
	\ite capstone display: 10 LEDs with \SI{9.8}{\milli\ampere} each; several switch at the same time
	\itee \SI{50}{\milli\ampere} through \SI{0.1}{\ohm} of ground: \SI{5}{\milli\volt}
	\itee at the bridge (\SI{4.65}{\milli\volt} per LSB) that would be one whole LSB!
	\ite at high frequencies the inductance dominates ($\approx\SI{1}{\nano\henry\per\milli\metre}$): \SI{10}{\nano\henry} and an edge of \SI{50}{\milli\ampere} in \SI{10}{\nano\second}: \SI{50}{\milli\volt} spike
\ei
\os{\newpage}

\item[{\rot\bf Common return impedance}]~\\[-\lblskip]
\begin{center}
\begin{circuitikz}[noifig=1.1]
  \draw (0,0) node[op amp, noinv input up](OA){};
  \coordinate (X) at ($(OA.out)+(0.6,0)$);
  \draw (OA.+) to[short, -o] ++(-1,0) node[left]{$\Vin$};
  \draw (OA.-) -- ++(0,-1.0) coordinate(N);
  \draw (N) to[R, l_=\SI{97.6}{\kilo\ohm}] (N -| X) -- (X);
  \draw (OA.out) -- (X) to[short,-o] ++(0.8,0) node[right]{$\Vout$};
  \draw (N) to[R, l_=\SI{10}{\kilo\ohm}] ++(0,-2.0) coordinate(G);
  \draw (G) to[R, l=\SI{0.1}{\ohm}, *-] ++(-3.0,0) coordinate(GG) -- ++(0,-0.2) node[ground]{};
  \draw (G) -- ++(6.0,0) coordinate(G2) to[R, l_=LEDs] ++(0,2.0) node[right]{\SI{5}{\volt}, switched};
  \draw[->, rwuviolet, thick] ($(G2)+(-0.3,-0.35)$) -- node[below, font=\scriptsize]{\SI{50}{\milli\ampere}} ($(GG)+(0.5,-0.35)$);
\end{circuitikz}
\end{center}
\bi
	\ite the gain stage's \SI{10}{\kilo\ohm} returns through the same trace as the LED current
	\ite its ground moves by \SI{5}{\milli\volt}: $\Vout$ changes by $-\frac{R_2}{R_1}\cdot\SI{5}{\milli\volt} = -\SI{49}{\milli\volt}$
	\ite Falstad: \SI{2.000}{\volt} (LEDs off) and \SI{1.951}{\volt} (on); star ground: always \SI{2.000}{\volt}
\ei
\falstad{noise_shared_ground}\par
\os{\newpage}

\item[{\rot\bf Remedies on the board}]~\\[-\lblskip]
\bi
	\ite \textbf{star ground}: every stage gets its own return to one common point (at the ADC or the reference)
	\ite separate \textbf{analog and digital returns} (AGND, DGND), joined at one point
	\ite a solid \textbf{ground plane}: very low resistance and inductance; return currents flow under their trace
	\ite large currents (LEDs, relays, motors) never through the ground of the input stage
	\ite \textbf{differential} inputs and a reference pin connected where the signal is used
	\ite numbers: \SI{49}{\milli\volt} is 0.1 LSB for 3 bit -- but 50 LSB for 12 bit, and it switches slowly enough to pass the anti-aliasing filter
\ei
\nsl{A low-pass filter does not help against a common-impedance error if
the disturbing current changes slowly. The LEDs of the capstone change
when the code changes, a few times per second or less: their ground drop
appears like a signal at the ADC. In the capstone the error is a tenth of
an LSB and does not matter, but the same layout with a 12-bit converter
would be useless. The ground connection must be designed, not left to the
autorouter.}
\os{\newpage}

\item[{\rot\bf Ground loops}]~\\[-\lblskip]
\begin{center}
\begin{circuitikz}[noifig=0.95]
  \draw (0,1.2) rectangle (2.4,2.8); \node[align=center, font=\small] at (1.2,2.0) {sensor\\device};
  \draw (8,1.2) rectangle (10.4,2.8); \node[align=center, font=\small] at (9.2,2.0) {amplifier\\(PC, logger)};
  \draw (2.4,2.4) -- node[above, font=\scriptsize]{signal} (8,2.4);
  \draw (2.4,1.6) to[R, l={$R_\mathrm{cable}$}] (8,1.6);
  \draw (1.2,1.2) -- (1.2,0) to[sV, l_=$\Delta V$] (4.2,0) to[R, l_=PE] (9.2,0) -- (9.2,1.2);
\end{circuitikz}
\end{center}
\bi
	\ite two devices, both connected to protective earth (PE), and a cable whose ground wire or shield ($R_\mathrm{cable}$) connects them again: a \textbf{loop}
	\ite earth points differ by \SI{10}{\milli\volt} to \SI{1}{\volt} (\SI{50}{\hertz}, harmonics): leakage currents of mains filters, induced voltage in the loop
	\ite the loop current flows through $R_\mathrm{cable}$; the drop adds to a \textbf{single-ended} signal
\ei
\os{\newpage}

\item[{\rot\bf Ground loop in Falstad}]~\\[-\lblskip]
\bi
	\ite \SI{1}{\volt} between the earth points, PE path \SI{4.5}{\ohm}, cable ground \SI{0.5}{\ohm}: \SI{200}{\milli\ampere}, remote ground moves by \SI{0.1}{\volt}
	\ite sensor signal \SI{37.2}{\milli\volt} (capstone full scale), amplifiers with $G = 10$:
\ei
\begin{center}
\begin{tabular}{lcc}
\toprule
input & DC output & \SI{50}{\hertz} at the output \\
\midrule
single-ended & \SI{0.372}{\volt} & \SI{1.0}{\volt} \\
difference amplifier, ideal resistors & \SI{0.372}{\volt} & 0 \\
difference amplifier, one resistor \SI{1}{\percent} off & \SI{0.372}{\volt} & \SI{0.90}{\milli\volt} \\
\bottomrule
\end{tabular}
\end{center}
\falstad{noise_ground_loop}\par
\bi
	\ite the difference amplifier measures signal minus remote ground: the ground difference is common mode
	\ite CMRR with a \SI{1}{\percent} resistor: $20\log(10/0.009) = \SI{61}{\deci\bel}$; integrated INA: \SIrange{90}{120}{\deci\bel}
\ei
\os{\newpage}

\item[{\rot\bf Remedies against ground loops}]~\\[-\lblskip]
\bi
	\ite \textbf{measure differentially}: run a separate sense wire to the remote ground, use an INA (as in the capstone)
	\ite \textbf{break the loop}
	\itee battery-powered or isolated device (isolated DC/DC converter and digital isolator, isolation amplifier, optocoupler)
	\itee shield grounded at \textbf{one end only} (low frequencies); for RF: other end via a \SI{10}{\nano\farad} capacitor
	\itee a ground-break resistor of \SIrange{10}{100}{\ohm} in the signal ground (the loop current drops; PE stays intact)
	\ite \textbf{never} remove the protective earth of a mains device -- it is a safety function
	\ite typical case in the lab: laptop with USB at the board, oscilloscope ground clip, sensor housing grounded at the machine
\ei
\os{\newpage}
\end{description}

%%%%%%%%%%%%%%%%%%%%%%%%%%%%%%%%%%%%%%%%%%%%%%%%%%%%%%%%%%%%%%%%%%%%%%%%%%%
\section{Supply and Reference Disturbances}
%%%%%%%%%%%%%%%%%%%%%%%%%%%%%%%%%%%%%%%%%%%%%%%%%%%%%%%%%%%%%%%%%%%%%%%%%%%
\nsl{The supply voltage carries ripple from the rectifier (\SI{100}{\hertz}),
switching noise from a regulator (\SI{100}{\kilo\hertz} to a few
\si{\mega\hertz}), and voltage drops caused by the load currents of other
circuits. An op-amp suppresses supply variations by its power supply
rejection ratio (PSRR), but the PSRR falls with frequency. A reference or
a bridge excitation passes its noise on directly. Decoupling, filtering
and the ratiometric principle keep these disturbances out.\newpage}

\begin{description}
\item[{\rot\bf Power supply rejection}]~\\[-\lblskip]
\bi
	\ite PSRR: supply change divided by the input-referred change it causes
	\itee typical \SIrange{80}{120}{\deci\bel} at DC, falling \SI{20}{\deci\bel} per decade: often only \SIrange{20}{40}{\deci\bel} at \SI{100}{\kilo\hertz}
	\ite example: \SI{50}{\milli\volt}\,pp ripple at \SI{100}{\kilo\hertz}, PSRR \SI{40}{\deci\bel}: \SI{0.5}{\milli\volt}\,pp at the input, $\times 107.6$
	\itee lies far above $f_c$ and is filtered -- if the stage stays linear (no clipping, no rectification)
	\ite stages \textbf{behind} the filter (S\&H, comparators, ladder) see their supply and reference without filter
	\ite remedies: \SI{100}{\nano\farad} at every supply pin (short!), \SIrange{1}{10}{\micro\farad} per board section, RC or LC filter, linear post-regulator (LDO)
\ei
\os{\newpage}

\item[{\rot\bf RC filter for the analog supply}]~\\[-\lblskip]
\begin{center}
\begin{circuitikz}[noifig=1.1]
  \draw (0,2) node[left]{\SI{5}{\volt} + ripple} to[short, o-] (0.5,2) to[R, l=\SI{10}{\ohm}] (2.8,2) -- (4.6,2) to[short, -o] (5.4,2) node[right]{analog \SI{5}{\volt}};
  \draw (2.8,2) to[C, l=\SI{10}{\micro\farad}, *-] (2.8,0) node[ground]{};
  \draw (4.6,2) to[R, l=\SI{1}{\milli\ampere} load, *-] (4.6,0) node[ground]{};
\end{circuitikz}
\end{center}
\bi
	\ite $f_c = 1/(2\pi\cdot\SI{10}{\ohm}\cdot\SI{10}{\micro\farad}) = \SI{1.6}{\kilo\hertz}$; at \SI{100}{\kilo\hertz}: $1/63$
	\ite Falstad: triangle ripple \SI{50}{\milli\volt}\,pp $\rightarrow$ \SI{0.63}{\milli\volt}\,pp; DC drop $\SI{1}{\milli\ampere}\cdot\SI{10}{\ohm} = \SI{10}{\milli\volt}$
	\ite the resistor must be small: the load current causes a drop that depends on the load
\ei
\falstad{noise_supply_filter}\par
\os{\newpage}

\item[{\rot\bf Reference noise and the ratiometric principle}]~\\[-\lblskip]
\bi
	\ite the bridge output is proportional to the excitation: \SI{1}{\percent} reference noise gives \SI{1}{\percent} signal noise
	\ite capstone: excitation and ADC ladder from the same \SI{4.00}{\volt}: noise and drift of the reference cancel (chapter~\ref{ch:References})
	\ite but only for disturbances that reach both paths \textbf{at the same time}
	\itee the signal path has a low-pass (\SI{68}{\milli\second}) and a S\&H, the ladder does not
	\itee slow drift cancels, fast reference noise does not -- filter the reference (RC after the reference, noise-reduction pin)
	\ite a voltage drop on a shared connection between bridge and ladder breaks the ratiometric principle
\ei
\falstad{ratiometric}\par
\os{\newpage}
\end{description}

%%%%%%%%%%%%%%%%%%%%%%%%%%%%%%%%%%%%%%%%%%%%%%%%%%%%%%%%%%%%%%%%%%%%%%%%%%%
\section{Thermoelectric Voltages and Other DC Errors}
%%%%%%%%%%%%%%%%%%%%%%%%%%%%%%%%%%%%%%%%%%%%%%%%%%%%%%%%%%%%%%%%%%%%%%%%%%%
\nsl{At DC and very low frequencies some disturbances look like a wrong
offset. They are not random, but they change with temperature and time,
so a single calibration does not remove them. The most important one in
low-level circuits is the thermoelectric voltage of junctions between
different metals. Others have been treated in earlier chapters and are
collected here for the error budget.\newpage}

\begin{description}
\item[{\rot\bf Thermoelectric voltages (Seebeck effect)}]~\\[-\lblskip]
\bi
	\ite every junction of two different metals is a thermocouple: $V = S\cdot\Delta T$
	\ite in a circuit the junctions come in pairs; they cancel only if both are at the \textbf{same temperature}
\ei
\begin{center}
\begin{tabular}{lc}
\toprule
junction & Seebeck coefficient \\
\midrule
copper -- copper & $<\SI{0.2}{\micro\volt\per\kelvin}$ \\
copper -- tin/lead solder & \SIrange{1}{5}{\micro\volt\per\kelvin} \\
copper -- Kovar (IC leads, relay contacts) & \SI{40}{\micro\volt\per\kelvin} \\
copper -- copper oxide (corroded contact) & \SI{1000}{\micro\volt\per\kelvin} \\
\bottomrule
\end{tabular}
\end{center}
\bi
	\ite a temperature difference of a few kelvin (warm regulator, airflow) gives several \si{\micro\volt} to \SI{100}{\micro\volt}
	\ite capstone: \SI{20}{\micro\volt} = \SI{0.05}{\kelvin}: nothing for 3 bit, but 2.2 LSB for 12 bit
\ei
\os{\newpage}

\item[{\rot\bf Remedies: symmetry and excitation reversal}]~\\[-\lblskip]
\bi
	\ite symmetric layout: the junctions of both signal wires side by side (same temperature, same materials)
	\ite keep heat sources (regulators, LED resistors) away from the input; avoid airflow over the input
	\ite \textbf{excitation reversal}: measure with $+V_B$ and $-V_B$
	\itee the bridge signal changes sign, the thermal EMF does not
\ei
\keyeq{Excitation reversal}{$V_1 = V_d + V_\mathrm{th}, \quad V_2 = -V_d + V_\mathrm{th} \qquad\Rightarrow\qquad V_d = \frac{V_1 - V_2}{2}, \quad V_\mathrm{th} = \frac{V_1 + V_2}{2}$}
\bi
	\ite Falstad: Pt100 at \SI{50}{\celsius}, $V_d = \SI{18.68}{\milli\volt}$, \SI{20}{\micro\volt} EMF: $V_1 = \SI{18.70}{\milli\volt}$, $V_2 = \SI{-18.66}{\milli\volt}$
\ei
\falstad{noise_thermal_emf}\par
\nsl{Excitation reversal also removes the offset and the offset drift of
the amplifier, and with a fast square-wave excitation it moves the signal
above the $1/f$ corner. Precision bridge amplifiers and many
delta-sigma converters for weigh scales use this technique (AC
excitation). It is the same idea as the chopper amplifier, applied to the
complete measurement path.}
\os{\newpage}

\item[{\rot\bf Other DC errors of the capstone (earlier chapters)}]~\\[-\lblskip]
\begin{center}
\begin{tabular}{lll}
\toprule
effect & size & remedy (chapter) \\
\midrule
op-amp offset & \SI{2}{\milli\volt} $\times$ 107.6 = \SI{0.22}{\volt} & zero-drift, trim (\ref{ch:OpAmps}) \\
offset drift & \SI{5}{\micro\volt\per\kelvin}: \SI{25}{\kelvin} $\rightarrow$ \SI{13}{\milli\volt} & zero-drift (\ref{ch:OpAmps}) \\
lead resistance (2-wire) & \SI{2}{}$\times$\SI{0.5}{\ohm} $\rightarrow$ \SI{2.6}{\kelvin} & 3-wire (\ref{ch:Bridges}) \\
self-heating & \SI{0.1}{\milli\watt}: \SIrange{0.01}{0.1}{\kelvin} & small current (\ref{ch:Bridges}) \\
reference drift & cancels & ratiometric (\ref{ch:References}) \\
\bottomrule
\end{tabular}
\end{center}
\bi
	\ite the offset of the first stage is the largest DC error of the capstone: almost half an LSB
\ei
\os{\newpage}
\end{description}

%%%%%%%%%%%%%%%%%%%%%%%%%%%%%%%%%%%%%%%%%%%%%%%%%%%%%%%%%%%%%%%%%%%%%%%%%%%
\section{Noise and Interference in the Sampled System}
%%%%%%%%%%%%%%%%%%%%%%%%%%%%%%%%%%%%%%%%%%%%%%%%%%%%%%%%%%%%%%%%%%%%%%%%%%%
\nsl{The last stages of the chain sample and quantize. Sampling folds
every frequency above $f_s/2$ into the band from 0 to $f_s/2$: noise and
interference that pass the anti-aliasing filter appear in the data as slow
fluctuations. On the other hand, a sampled system can average, which
reduces random noise.\newpage}

\begin{description}
\item[{\rot\bf Aliasing of noise and interference}]~\\[-\lblskip]
\bi
	\ite \SI{50}{\hertz} sampled at \SI{5}{\hertz} ($50 = 10\cdot f_s$): aliases to \SI{0}{\hertz} -- a constant error of unknown size (chapters~\ref{ch:Filters}, \ref{ch:ADC})
	\ite white noise: all noise up to the bandwidth of the amplifier folds into $0 \ldots f_s/2$
	\itee without filter: noise bandwidth some \SI{10}{\kilo\hertz} instead of \SI{3.68}{\hertz}: $\sqrt{10000/3.68} = 50$ times more noise in the samples
	\ite the anti-aliasing filter limits the noise bandwidth as well -- one more reason to place it \textbf{before} the S\&H
\ei
\falstad{alias50}\quad\falstad{capstone_nofilter}\par
\os{\newpage}

\item[{\rot\bf Averaging}]~\\[-\lblskip]
\bi
	\ite the mean of $N$ samples with independent noise: noise RMS $/\sqrt{N}$
	\itee 16 samples: noise $/4$; 100 samples: noise $/10$ -- at the cost of $N$ times the measurement time
	\ite works only for noise that is \textbf{independent} from sample to sample (white noise above $f_s$)
	\itee not for $1/f$ noise and drift; not for interference synchronous with $f_s$ (\SI{50}{\hertz} at \SI{5}{\hertz})
	\ite \textbf{integration over whole mains periods} (\SI{20}{\milli\second}, \SI{100}{\milli\second}): \SI{50}{\hertz} and all harmonics average to zero -- dual-slope and delta-sigma converters (chapter~\ref{ch:ADC})
	\ite a little noise (dither, about 0.5 LSB RMS) before an ADC allows averaging to resolve below one LSB
\ei
\os{\newpage}

\item[{\rot\bf Noise at a comparator threshold}]~\\[-\lblskip]
\bi
	\ite a slow input near a threshold plus noise: the comparator output chatters
	\ite hysteresis larger than the peak-to-peak noise (chapter~\ref{ch:ADC})
	\ite in the flash converter: noise can produce bubbles in the thermometer code
	\ite the S\&H holds the input still during the decision; the anti-aliasing filter keeps the held value quiet
\ei
\falstad{comparator_hyst}\par
\os{\newpage}
\end{description}

%%%%%%%%%%%%%%%%%%%%%%%%%%%%%%%%%%%%%%%%%%%%%%%%%%%%%%%%%%%%%%%%%%%%%%%%%%%
\section{Error Budget and Design Rules}
%%%%%%%%%%%%%%%%%%%%%%%%%%%%%%%%%%%%%%%%%%%%%%%%%%%%%%%%%%%%%%%%%%%%%%%%%%%
\nsl{We now put all numbers of this chapter together for the capstone
chain and for the same chain with a 12-bit converter. Random noise is given
as peak-to-peak value, interference as amplitude, both at the ADC input
(after the low-pass filter, where it applies).\newpage}

\begin{description}
\item[{\rot\bf Disturbances of the capstone at the ADC input}]~\\[-\lblskip]
\begin{center}
\begin{tabular}{lccc}
\toprule
disturbance & at the ADC & 3 bit & 12 bit \\
 & & (\SI{0.5}{\volt}) & (\SI{0.98}{\milli\volt}) \\
\midrule
op-amp and resistor noise (pp) & \SI{0.14}{\milli\volt} & 0.0003 & 0.14 \\
\SI{50}{\hertz} magnetic, \SI{5}{\milli\volt}, filtered & \SI{25}{\milli\volt} & 0.05 & 25 \\
\SI{50}{\hertz} capacitive, \SI{10}{\micro\volt} at the tap, filtered & \SI{0.05}{\milli\volt} & 0.0001 & 0.05 \\
LED current in the ground of $G_2$ & \SI{49}{\milli\volt} & 0.1 & 50 \\
thermal EMF \SI{20}{\micro\volt} & \SI{2.2}{\milli\volt} & 0.004 & 2.2 \\
op-amp offset \SI{2}{\milli\volt} (untrimmed) & \SI{0.22}{\volt} & 0.43 & 220 \\
quantization noise (RMS) & -- & 0.29 & 0.29 \\
\bottomrule
\end{tabular}
\end{center}
\bi
	\ite 3 bit: only the offset needs attention (trim or zero-drift op-amp) -- the capstone design is sufficient
	\ite 12 bit: offset, layout (ground), \SI{50}{\hertz} (second-order filter, twisted pair) and thermal EMF must all be handled
\ei
\os{\newpage}

\item[{\rot\bf Design card}]~\\[-\lblskip]
\designcard{Low-noise, interference-free measurement front end}{
\begin{enumerate}
\item \textbf{Yardstick}: calculate the LSB at every point of the chain (capstone: \SI{4.65}{\milli\volt} at the bridge, \SI{0.5}{\volt} at the ADC). Aim for each disturbance $< \frac14$ LSB.
\item \textbf{Bandwidth}: no more than needed. The anti-aliasing filter before the S\&H sets the noise bandwidth ($B_n = 1.57 f_c$).
\item \textbf{First stage}: gain $\geq 10$, low $e_n$ for low-ohmic sources, low $i_n$ for high-ohmic ones; zero-drift for DC; metal-film resistors.
\item \textbf{Sensor cable}: twisted pair, shield grounded at the amplifier only, differential input (INA), RC filter against RF at the input.
\item \textbf{Ground}: star point or plane; LED, relay and digital currents never through the input ground; no ground loops (one earth connection or isolation).
\item \textbf{Supply and reference}: \SI{100}{\nano\farad} at every IC, RC/LC filter for the analog part, ratiometric reference.
\item \textbf{DC}: symmetric junctions, heat sources away from the input; offset trim or excitation reversal.
\end{enumerate}}
\os{\newpage}
\end{description}

%%%%%%%%%%%%%%%%%%%%%%%%%%%%%%%%%%%%%%%%%%%%%%%%%%%%%%%%%%%%%%%%%%%%%%%%%%%
\section{Exercises}
%%%%%%%%%%%%%%%%%%%%%%%%%%%%%%%%%%%%%%%%%%%%%%%%%%%%%%%%%%%%%%%%%%%%%%%%%%%

\exercise{Thermal noise and bandwidth}\label{ex:noise-thermal}
\begin{talist}
\ta Calculate the noise density of a \SI{10}{\kilo\ohm} resistor at \SI{300}{\kelvin} and its RMS noise voltage in a bandwidth of \SI{10}{\kilo\hertz}.
\ta The resistor is followed by an RC low-pass with $f_c = \SI{100}{\hertz}$. Calculate the noise bandwidth and the RMS noise voltage behind the filter. Estimate the peak-to-peak value.
\ta Why is the result of b) independent of the capacitor's own noise?
\ta Falstad's noise generator with amplitude \SI{1}{\volt} and time step \SI{10}{\micro\second} has the density \SI{2.58}{\milli\volt}\rtHz. Calculate the RMS values behind low-passes with \SI{995}{\hertz} and \SI{99.5}{\hertz} and check them in Falstad.
\end{talist}

\begin{solution}
\begin{talist}
\ta $e_n = \sqrt{\SI{1.66e-20}{\joule}\cdot\SI{10}{\kilo\ohm}} = \SI{12.9}{\nano\volt}\rtHz$;
	$V_\mathrm{rms} = \SI{12.9}{\nano\volt}\rtHz\cdot\sqrt{\SI{10}{\kilo\hertz}} = \SI{1.29}{\micro\volt}$.
\ta $B_n = \frac{\pi}{2}\cdot\SI{100}{\hertz} = \SI{157}{\hertz}$;
	$V_\mathrm{rms} = \SI{12.9}{\nano\volt}\rtHz\cdot\sqrt{\SI{157}{\hertz}} = \SI{0.16}{\micro\volt}$;
	$V_\mathrm{pp} \approx 6.6\cdot\SI{0.16}{\micro\volt} = \SI{1.1}{\micro\volt}$.
\ta An ideal capacitor has no resistance (no real part of its impedance) and therefore no thermal noise.
\ta $B_n = 1.57\cdot\SI{995}{\hertz} = \SI{1563}{\hertz}$: $\SI{2.58}{\milli\volt}\rtHz\cdot\sqrt{1563} = \SI{0.102}{\volt}$;
	$B_n = \SI{156}{\hertz}$: \SI{0.032}{\volt}.
	\falstad{noise_bandwidth}: the scopes show about \SI{0.10}{\volt} and \SI{0.032}{\volt} RMS
	(a few percent different on every run), the unfiltered generator \SI{0.577}{\volt}.
\end{talist}
\end{solution}

\exercise{Noise of the capstone front end}\label{ex:noise-ina}
The INA of the capstone consists of two buffers and a difference amplifier
(\SI{10}{\kilo\ohm}/\SI{100}{\kilo\ohm}, $G_1 = 10$), followed by the gain stage
$G_2 = 10.76$ and the low-pass with $f_c = \SI{2.34}{\hertz}$.
\begin{talist}
\ta Calculate the thermal noise density of the bridge as seen by the INA (both taps, Pt100 at \SI{0}{\celsius}).
\ta The op-amps have $e_n = \SI{55}{\nano\volt}\rtHz$ and no $1/f$ noise. The difference amplifier op-amp is amplified with the noise gain 11. Calculate the input-referred noise density of the INA without resistor noise.
\ta Calculate the RMS noise at the ADC input and its peak-to-peak value. Compare with the LSB of a 3-bit and a 12-bit converter.
\ta Repeat c) for op-amps with $e_w = \SI{40}{\nano\volt}\rtHz$ and $f_{nc} = \SI{100}{\hertz}$, for a measurement time of \SI{100}{\second}.
\end{talist}

\begin{solution}
\begin{talist}
\ta Each tap: $\SI{3.9}{\kilo\ohm}\parallelto\SI{100}{\ohm} = \SI{97.5}{\ohm}$; together \SI{195}{\ohm}:
	$\sqrt{\SI{1.66e-20}{\joule}\cdot\SI{195}{\ohm}} = \SI{1.8}{\nano\volt}\rtHz$.
\ta Buffers: $\sqrt{2}\cdot\SI{55}{\nano\volt}\rtHz = \SI{78}{\nano\volt}\rtHz$. Difference amplifier: its op-amp noise appears
	at the output with gain 11, referred to the input (gain 10): $1.1\cdot\SI{55}{\nano\volt}\rtHz = \SI{61}{\nano\volt}\rtHz$.
	RSS: $\sqrt{78^2 + 61^2}\,\si{\nano\volt}\rtHz = \SI{99}{\nano\volt}\rtHz$ (with the bridge and the
	resistors, \SI{19}{\nano\volt}\rtHz: \SI{100}{\nano\volt}\rtHz).
\ta $B_n = 1.57\cdot\SI{2.34}{\hertz} = \SI{3.68}{\hertz}$; input: $\SI{100}{\nano\volt}\rtHz\cdot\sqrt{3.68} = \SI{0.19}{\micro\volt}$;
	ADC input: $\times 107.6 = \SI{21}{\micro\volt}$ RMS, $\approx\SI{0.14}{\milli\volt}$\,pp.
	3 bit: $\SI{0.14}{\milli\volt}/\SI{0.5}{\volt} = 0.0003$ LSB; 12 bit: $0.14/0.98 = 0.14$ LSB. Uncritical.
\ta $f_1 = \SI{0.01}{\hertz}$: per op-amp $\SI{40}{\nano\volt}\rtHz\cdot\sqrt{100\ln(368) + 3.67}\,\sqrt{\mathrm{Hz}} = \SI{0.98}{\micro\volt}$.
	Three op-amps (factors 1, 1, 1.1): $\SI{0.98}{\micro\volt}\cdot\sqrt{1+1+1.21} = \SI{1.75}{\micro\volt}$;
	output: \SI{0.19}{\milli\volt} RMS, \SI{1.2}{\milli\volt}\,pp: about one LSB of a 12-bit converter, still nothing for 3 bit.
\end{talist}
\end{solution}

\exercise{Capacitive coupling}\label{ex:noise-cap}
A cable carrying \SI{230}{\volt}, \SI{50}{\hertz} runs next to a signal line; the stray capacitance is \SI{1}{\pico\farad}.
\begin{talist}
\ta Calculate the amplitude of the coupled current.
\ta Calculate the interference amplitude on the signal line if the line ends in \SI{1}{\mega\ohm}, \SI{10}{\kilo\ohm} or at the capstone bridge tap (\SI{97.5}{\ohm}).
\ta The \SI{1}{\mega\ohm} line is shielded; \SI{1}{\percent} of the capacitance remains. What is left?
\ta What happens if the shield is not connected?
\ta Check b) and c) in Falstad.
\end{talist}

\begin{solution}
\begin{talist}
\ta $\hat\imath = \omega C\hat V = 2\pi\cdot\SI{50}{\hertz}\cdot\SI{1}{\pico\farad}\cdot\SI{325}{\volt} = \SI{0.102}{\micro\ampere}$.
\ta $1/(\omega C) = \SI{3.2}{\giga\ohm}$ is much larger than every node impedance, so $\hat v = \hat\imath R$:
	\SI{0.10}{\volt}, \SI{1.0}{\milli\volt} and \SI{10}{\micro\volt}. At the bridge: $\times 107.6 = \SI{1.1}{\milli\volt}$ before
	the filter (and less if it couples into both taps as common mode).
\ta \SI{0.01}{\pico\farad}: \SI{1.0}{\milli\volt}.
\ta A floating shield takes the potential of its surroundings and couples to the line with its own (large) capacitance:
	it may even make things worse. The shield must be connected to the signal reference.
\ta \falstad{noise_cap_coupling}\newline \texttt{V\_1M} \SI{0.102}{\volt}, \texttt{V\_10k} \SI{1.02}{\milli\volt},
	\texttt{V\_100} \SI{10.2}{\micro\volt}, \texttt{V\_shield} \SI{1.02}{\milli\volt} (amplitudes).
\end{talist}
\end{solution}

\exercise{Ground loop and difference amplifier}\label{ex:noise-loop}
A remote sensor delivers \SI{37.2}{\milli\volt} against its own ground. Between the earth points of the
sensor and the amplifier there are \SI{1}{\volt} at \SI{50}{\hertz}; the PE path has \SI{4.5}{\ohm}, the ground
conductor of the cable \SI{0.5}{\ohm}.
\begin{talist}
\ta Calculate the loop current and the voltage between the two grounds.
\ta The signal is amplified single-ended with $G=10$. What is at the output?
\ta A difference amplifier (\SI{10}{\kilo\ohm}/\SI{100}{\kilo\ohm}) measures between the signal and a sense wire to the remote ground.
	One of the \SI{100}{\kilo\ohm} resistors is \SI{1}{\percent} too large. Calculate the output and the CMRR.
\ta Name two ways to break the loop.
\ta Check b) and c) in Falstad.
\end{talist}

\begin{solution}
\begin{talist}
\ta $I = \SI{1}{\volt}/(\SI{4.5}{\ohm} + \SI{0.5}{\ohm}) = \SI{0.2}{\ampere}$; $\Delta V = \SI{0.2}{\ampere}\cdot\SI{0.5}{\ohm} = \SI{0.1}{\volt}$ amplitude.
\ta $10\cdot(\SI{37.2}{\milli\volt} + \SI{0.1}{\volt}\sin\omega t)$: \SI{0.372}{\volt} DC with \SI{1.0}{\volt} of \SI{50}{\hertz} -- the hum is almost three times the signal.
\ta With $V^+$ side \SI{10}{\kilo\ohm}/\SI{101}{\kilo\ohm}: $\Vout = 11\cdot\frac{101}{111}(V_g + V_s) - 10\,V_g = 10.009\,V_s + 0.009\,V_g$.
	Output \SI{0.372}{\volt} DC with \SI{0.90}{\milli\volt} of \SI{50}{\hertz}. CMRR $= 20\log(10/0.009) = \SI{61}{\deci\bel}$.
\ta Battery or isolated supply for one device (isolated DC/DC converter, digital isolator, isolation amplifier); ground the
	cable shield at one end only; ground-break resistor in the signal ground. The protective earth must stay.
\ta \falstad{noise_ground_loop}\newline \texttt{GND\_remote} \SI{0.100}{\volt} amplitude; \texttt{Vout\_se} \SI{0.372}{\volt} $\pm$ \SI{1.00}{\volt};
	\texttt{Vout\_diff} \SI{0.372}{\volt} without hum; \texttt{Vout\_diff1} \SI{0.372}{\volt} $\pm$ \SI{0.90}{\milli\volt}.
\end{talist}
\end{solution}

\exercise{Shared ground and thermal EMF}\label{ex:noise-ground}
\begin{talist}
\ta The gain stage of the capstone ($G_2 = 1 + \SI{97.6}{\kilo\ohm}/\SI{10}{\kilo\ohm}$) returns its \SI{10}{\kilo\ohm} resistor through
	\SI{0.1}{\ohm} of ground trace that also carries \SI{50}{\milli\ampere} of LED current. Calculate the error at the output and express it in LSB for 3 and 12 bit.
\ta Why does the anti-aliasing filter not remove this error?
\ta A thermal EMF of \SI{20}{\micro\volt} is in series with the bridge signal. Express it in kelvin and in LSB (3 bit, 12 bit).
\ta The bridge is measured with $+\SI{4}{\volt}$ and $-\SI{4}{\volt}$ excitation at \SI{50}{\celsius} ($R = \SI{119.25}{\ohm}$). Calculate both readings and how to get $V_d$ and the EMF from them.
\ta Check a) and d) in Falstad.
\end{talist}

\begin{solution}
\begin{talist}
\ta Ground shift $\SI{50}{\milli\ampere}\cdot\SI{0.1}{\ohm} = \SI{5}{\milli\volt}$; it acts on the inverting side with gain $-R_2/R_1 = -9.76$:
	\SI{-49}{\milli\volt}. 3 bit: 0.1 LSB; 12 bit: 50 LSB.
\ta The LED current changes only when the code changes, i.e.\ slowly, within the pass band of the filter. Only a separate return (star ground) helps.
\ta $\SI{20}{\micro\volt}/(\SI{0.372}{\milli\volt\per\kelvin}) = \SI{0.054}{\kelvin}$; at the ADC $\times 107.6 = \SI{2.2}{\milli\volt}$: 0.004 LSB (3 bit), 2.2 LSB (12 bit).
\ta $V_d = \SI{4}{\volt}\cdot(119.25/4019.25 - 100/4000) = \SI{18.68}{\milli\volt}$; $V_1 = \SI{18.70}{\milli\volt}$, $V_2 = \SI{-18.66}{\milli\volt}$;
	$(V_1 - V_2)/2 = \SI{18.68}{\milli\volt}$, $(V_1 + V_2)/2 = \SI{20}{\micro\volt}$.
\ta \falstad{noise_shared_ground}\newline \texttt{Vout\_shared} \SI{2.000}{\volt} / \SI{1.951}{\volt}, \texttt{Vout\_star} \SI{2.000}{\volt};\newline
	\falstad{noise_thermal_emf}\newline voltmeters \SI{18.70}{\milli\volt} and \SI{-18.66}{\milli\volt}, INA outputs \texttt{V1x10} \SI{0.1870}{\volt}, \texttt{V2x10} \SI{-0.1866}{\volt}.
\end{talist}
\end{solution}

%%% Local Variables:
%%% mode: latex
%%% TeX-master: "docu"
%%% End:
